\clearpage
\begin{examplebox}
  \begin{minipage}{\linewidth}
\textbf{\textcolor{CExample}{例 1：为什么有时端点必须去掉？（入门）}}
\begin{center}\begin{tikzpicture}[font=\small,>=stealth,line join=round]
\draw[gray!45,rounded corners=3pt] (0,0) rectangle (14.2,9.4);
\node at (7.1,8.8) {函数无限接近下边界，但从不落到下边界上};

\begin{scope}[shift={(4.1,1.8)},x=.9cm,y=1.65cm]
\draw[->] (-3.5,0)--(3.6,0) node[right]{$x$};\draw[->] (0,-.1)--(0,3.8) node[above]{$y$};\draw[dashed,magenta] (-3.3,2)--(3.3,2);
\draw[teal!80!black,thick,domain=-3.2:3.2,samples=100] plot(\x,{2+1/(\x*\x+1)});\fill[blue] (0,3) circle(2pt);\node[above right] at (0,3){$(0,3)$};\node[above left] at (0,2){$2$};
\end{scope}
\node[align=left,anchor=west] at (8.3,5.1){$y=2+\dfrac1{x^2+1}$\\[10pt] $0<\dfrac1{x^2+1}\le1$\\[10pt] 值域 $(2,3]$\\[12pt] $y=2$ 时变成 $1=0$};
\node[text width=13.2cm,align=center] at (7.1,.65) {即使候选判别式在 $y=2$ 处为零，降次后也可能无解。};
\end{tikzpicture}
\end{center}
\end{minipage}



  \medskip
  求 $y=\dfrac{2x^2+3}{x^2+1}$ 的值域。

  \medskip
  \textbf{第 1 步：}分母 $x^2+1>0$，所有实数 $x$ 都可以代入。

  \textbf{第 2 步：}当 $y\ne2$ 时，移项得
  \[
    (2-y)x^2+(3-y)=0.
  \]
  它要有实根，需
  \[
    \Delta(y)=-4(2-y)(3-y)\ge0
    \quad\Longrightarrow\quad 2\le y\le3.
  \]

  \textbf{第 3 步：}刚才用判别式时要求 $y\ne2$，所以再查 $y=2$：
  \[
    (2-2)x^2+(3-2)=0
    \quad\Longrightarrow\quad 1=0,
  \]
  不成立；而 $x=0$ 时 $y=3$，右端点能取到。

  \textbf{答案：}\quad $\boxed{y\in(2,3]}$。

  \medskip
  \textbf{\textcolor{CExample}{例 2：分子分母都有一次项（标准题）}}

  \medskip
  求 $y=\dfrac{x^2+1}{x^2+2x+2}$ 的值域。

  \medskip
  \textbf{第 1 步：}分母 $(x+1)^2+1>0$。

  \textbf{第 2 步：}去分母：
  \[
    (1-y)x^2-2yx+(1-2y)=0.
  \]
  对 $y\ne1$，要求
  \[
    \begin{aligned}
      \Delta(y)&=(-2y)^2-4(1-y)(1-2y)\\
      &=4(-y^2+3y-1)\ge0.
    \end{aligned}
  \]
  因此 $y^2-3y+1\le0$，解出
  \[
    \frac{3-\sqrt5}{2}\le y\le\frac{3+\sqrt5}{2}.
  \]

  \textbf{第 3 步：}检查 $y=1$：此时 $-2x-1=0$，解得 $x=-\dfrac12$，确实有解，因此 $1$ 也在值域中。

  \textbf{答案：}\quad $\boxed{y\in\left[\dfrac{3-\sqrt5}{2},\dfrac{3+\sqrt5}{2}\right]}$。

  \medskip
  \textbf{\textcolor{CExample}{例 3：二次项消失以后，反而变得更简单（提高）}}

  \medskip
  求 $y=\dfrac{3x^2+2x+1}{x^2+x+1}$ 的值域。

  \medskip
  \textbf{第 1 步：}分母
  \[
    x^2+x+1=\left(x+\frac12\right)^2+\frac34>0.
  \]
  \textbf{第 2 步：}去分母得
  \[
    (3-y)x^2+(2-y)x+(1-y)=0.
  \]
  对 $y\ne3$，要求
  \[
    \begin{aligned}
      \Delta(y)&=(2-y)^2-4(3-y)(1-y)\\
      &=-3y^2+12y-8\ge0.
    \end{aligned}
  \]
  整理得 $(y-2)^2\le\dfrac43$，所以
  \[
    2-\frac{2\sqrt3}{3}\le y\le2+\frac{2\sqrt3}{3}.
  \]
  \textbf{第 3 步：}检查 $y=3$：方程变为 $-x-2=0$，得 $x=-2$，说明 $y=3$ 能取到，而且它在上述区间内。

  \textbf{答案：}\quad $\boxed{y\in\left[2-\dfrac{2\sqrt3}{3},\,2+\dfrac{2\sqrt3}{3}\right]}$。

  \medskip
  \textbf{\textcolor{CExample}{例 4：所有 $x$ 都得到同一个 $y$（容易忽略）}}

  \medskip
  求 $y=\dfrac{2x^2+2}{x^2+1}$ 的值域。

  \medskip
  因为 $2x^2+2=2(x^2+1)$ 且 $x^2+1\ne0$，所以
  \[
    y=2.
  \]
  若坚持去分母，也会得到 $(2-y)(x^2+1)=0$；只有 $y=2$ 时方程成立。

  \textbf{答案：}\quad $\boxed{\{2\}}$。
  \par\textbf{提醒：}这是“$x^2$ 项消失，剩下的式子又恰好等于 $0$”的一种情况，不能漏掉。

  \medskip
  \textbf{\textcolor{CExample}{一句话比较}}
  例 1 的 $\dfrac ad=2$ \textbf{取不到}；例 2 的 $\dfrac ad=1$ 和例 3 的 $\dfrac ad=3$ \textbf{都取得到}；例 4 的函数\textbf{始终等于} $\dfrac ad=2$。这四题刚好覆盖了最容易混淆的情况。
\end{examplebox}
